\section{Parallel \cpi{} Construction}\label{sec:parabuild}

\peelH{} reduces residual support maintenance to static path-counter updates.
%
After this change, construction can become the main cost.
%
The phase breakdown in Section~\ref{sec:exp} reports cases where construction is a large fraction of total wall-clock.
%
This makes construction part of the algorithmic pipeline rather than a secondary implementation detail.

\stitle{Why the seed loop can be parallel.}
The \cpi{} builder follows the standard degeneracy-ordered Bron--Kerbosch search with Tomita pivoting~\cite{NuclearCD,BronKerbosch,eppstein,tomita}.
%
Each \(s\)-clique has one owner: the smallest vertex of the clique in the degeneracy order.
%
Therefore, the outer loop over seed vertices partitions the emitted \(s\)-clique families.
%
No two seeds own the same \(s\)-clique.

Static CPI also changes what construction must produce.
%
The mutable baseline builds residual tree state because later peeling may edit that state.
%
Our \(r=1\) pipeline never edits the stored hold or pivot sets.
%
The builder therefore only has to emit a path list and vertex--path incidence arrays.
%
This output can be accumulated in thread-local arenas and merged after all seeds finish.

\stitle{Algorithm state.}
\paraBuild{} takes \(G\), \(s\), and a thread count.
%
Each worker owns a contiguous seed range, a local path arena, and a generation-stamped membership bitmap.
%
The local invariant is that the worker arena contains exactly the paths emitted by seeds already processed by that worker.
%
The global merge concatenates the arenas in deterministic worker order and then builds the shared incidence CSR.
%
The cost driver is the Bron--Kerbosch search work under each seed, plus the linear merge over emitted incidences.

\begin{algorithm}[t]
\algfontsize
\caption{\paraBuild}
\label{alg:parabuild}
\KwIn{Graph \(G\) in CSR; clique size \(s\); thread count \(T\).}
\KwOut{Static \cpi{} \(\mathcal{P}\) in shared CSR.}

compute degeneracy ordering \(\sigma\) on \(G\)\;
launch \(T\) threads and assign each thread \(t\) a contiguous chunk \(\sigma_t\) of seed vertices\;
\ForEach{thread \(t\in\{1,\ldots,T\}\) \emph{in parallel}}{
    \(A_t\gets\emptyset\) \tcp*{thread-local path arena}
    \(b_t\gets\) generation-stamped bitmap over \(V\)\;
    \ForEach{seed vertex \(v\in\sigma_t\)}{
        increment generation of \(b_t\)\;
        \(N^+\gets\{u\in N_G(v)\mid \sigma(u)>\sigma(v)\}\)\;
        \textsc{BkEmit}\((v,N^+,A_t,b_t)\) \tcp*{pivoted BK on \(G[N^+]\), with \(v\) as a hold}
    }
}
\(\mathcal{P}\gets\) deterministic concat-and-prefix-sum merge of \(\{A_t\}_{t=1}^T\)\;
build incidence CSR over \(\mathcal{P}\) in parallel\;
\Return{\(\mathcal{P}\)}
\end{algorithm}

\textsc{BkEmit} is the same pivot-aware enumeration used by the CPI construction lineage~\cite{Pivoter,NuclearCD}.
%
The difference is the output contract.
%
It streams each emitted path into \(A_t\) instead of storing a mutable recursion tree for later residual updates.

\stitle{Correctness.}
\paraBuild{} is a parallel schedule of the same seed-owned enumeration.
%
Seed ownership gives disjoint emitted \(s\)-clique families.
%
The deterministic merge changes only the storage order, not the represented path families.
%
The resulting path set satisfies Definition~\ref{def:path}.
%
Therefore, the support formula of Lemma~\ref{lem:supp} applies unchanged.

\stitle{Scope.}
The design parallelizes construction because construction is the phase exposed by the static-\cpi{} pipeline.
%
It does not claim a lower bound showing that \peelH{} cannot be parallelized.
%
It also does not claim ideal scaling.
%
The scaling evidence is reported in Table~\ref{tab:par}.
% TODO: verify Table~\ref{tab:par} against the raw parallel-construction source file referenced by the plotting scripts.
